\section{An unconditional central-order certificate and exactly two right branches}
\label{sec:analytic-central-order-certificate}

The next argument uses only the classical critical strip, not a
Riemann hypothesis. It complements NN-A018 by determining the
central multiplicity of its explicit quaternionic example.
The central order itself is treated as a classical type of
calculation, without a historical novelty claim. The new
catalogue consequence is the complete two-branch conclusion for
the specified shifted deformation, together with the fully
written rational certificate below.

For a number field $K$ of signature $(r_1,r_2)$, use the exact
pole-removed completion
\begin{equation}\label{eq:analytic-order-completion}
 \xi_K(s)=s(s-1)D_K^{s/2}\pi^{-r_1s/2}(2\pi)^{-r_2s}
        \Gamma(s/2)^{r_1}\Gamma(s)^{r_2}\zeta_K(s).
\end{equation}
It is real entire of order one, is invariant under $s\mapsto1-s$,
and its zeros are the nontrivial zeros in $0\le\Re s\le1$,
with multiplicity. A primary locator for these analytic inputs
and the completion is Hasanalizade--Shen--Wong,
\emph{Counting zeros of Dedekind zeta functions},
\emph{Math. Comp.}~91 (2022), section~2.1, equations~(2.1)--(2.2),
p.~279. Its complex-place gamma convention differs from
\eqref{eq:analytic-order-completion} by the positive constant
$2^{r_2}$, by the duplication formula; every logarithmic
derivative here is therefore unchanged. The Euler product gives
nonvanishing in $\Re s>1$, and the functional equation and the
nonzero residues at $0,1$ remove any additional completed zeros.

\begin{theorem}[NN-A033: a positive rational test for central order]
\label{thm:analytic-central-order-test}
Put $F_K(z)=\xi_K(1/2+z)$,
$G_K(z)=F_K'(z)/F_K(z)$, and
$m_K=\operatorname{ord}_{s=1/2}\zeta_K(s)$. For every real
$\delta>1/\sqrt2$,
\begin{equation}\label{eq:analytic-T2-order}
 m_K\le T_{2,K}(\delta):=
 \frac{\delta G_K(\delta)-\delta^2G_K'(\delta)}2
 =m_K+\sum_{\{\alpha,-\alpha\}}
                 \frac{2\delta^4}{(\delta^2-\alpha^2)^2}.
\end{equation}
Here $\alpha=\rho-1/2\ne0$, one representative is chosen for
each pair $\{\alpha,-\alpha\}$, multiplicities are retained,
and conjugate terms are grouped. Each grouped contribution is
strictly positive. The series is absolutely convergent.

Set $\sigma=1/2+\delta$ and
\[
 B_K(s)=\frac1s+\frac1{s-1}+\frac12\log D_K
       -\frac{r_1}2\log\pi-r_2\log(2\pi)
       +\frac{r_1}2\psi(s/2)+r_2\psi(s).
\]
Then the entirely archimedean upper bound is
\begin{equation}\label{eq:analytic-T2-arch}
 T_{2,K}(\delta)<
 U_K(\delta):=
       \frac{\delta B_K(\sigma)-\delta^2B_K'(\sigma)}2.
\end{equation}
Both assertions are unconditional.
\end{theorem}

\begin{proof}
Evenness gives $m_K$ even. Hadamard factorization and pairing
opposite zeros give
\[
 F_K(z)=c z^{m_K}\prod_{\{\alpha,-\alpha\}}
                      (1-z^2/\alpha^2),\qquad c\ne0.
\]
The product converges normally because the order-one zero
divisor satisfies $\sum_{\alpha\ne0}|\alpha|^{-2}<\infty$.
Pairing the genus-one factors cancels their exponential parts;
any remaining factor $e^{az+b}$ has $a=0$ by evenness.
The point $\delta$ is not a zero, since
$\sigma>1$. Logarithmic differentiation on a neighbourhood of
that point is consequently legitimate and gives
\[
 G_K(z)=\frac{m_K}{z}
       +\sum_{\{\alpha,-\alpha\}}
                        \frac{2z}{z^2-\alpha^2}.
\]
Differentiating once more, the contribution of a noncentral
pair to $(\delta G_K-\delta^2G_K')/2$ is
$2\delta^4/(\delta^2-\alpha^2)^2$, and the central contribution
is $m_K$. This also proves absolute convergence of the series
in \eqref{eq:analytic-T2-order}, whose tail is
$O_\delta(\sum |\alpha|^{-4})$.

It remains to check positivity without moving zeros to the
critical line. Write $\alpha=a+ib$, where $|a|\le1/2$. Put
\[
 A=\delta^2-a^2+b^2,\qquad B=2ab.
\]
Then
\[
 A-|B|=\delta^2-2a^2+(|b|-|a|)^2
                    >0.
\]
For a genuine quartet $\alpha,-\alpha,\bar\alpha,-\bar\alpha$,
the two paired contributions therefore have positive real sum,
because
\[
 \Re\frac1{(\delta^2-\alpha^2)^2}
             =\frac{A^2-B^2}{(A^2+B^2)^2}>0.
\]
If the pair is real or purely imaginary, its single contribution
is itself positive. This proves the first inequality with all
possible off-line zeros included.

For real $\sigma>1$, the ideal Euler product gives the absolutely
convergent positive sums
\[
 E_1(\sigma)=\sum_{\mathfrak p}\sum_{k\ge1}
           \frac{\log N\mathfrak p}{(N\mathfrak p)^{k\sigma}}>0,
 \qquad
 E_2(\sigma)=\sum_{\mathfrak p}\sum_{k\ge1}
           \frac{k(\log N\mathfrak p)^2}
                {(N\mathfrak p)^{k\sigma}}>0.
\]
Thus $\zeta_K'/\zeta_K=-E_1$ and
$(\zeta_K'/\zeta_K)'=E_2$. Taking the logarithmic derivatives of
\eqref{eq:analytic-order-completion} now yields exactly
\[
 T_{2,K}(\delta)=U_K(\delta)
                  -\frac\delta2E_1(\sigma)
                  -\frac{\delta^2}2E_2(\sigma)
                  <U_K(\delta),
\]
which proves \eqref{eq:analytic-T2-arch}.
\end{proof}

\begin{corollary}[The quaternionic example has central order exactly two]
\label{cor:analytic-martinet-exact-two}
Let
\[
 F=\mathbb Q(\sqrt5,\sqrt{41}),\qquad
 N=F\left(\sqrt{\frac{(5+\sqrt5)(41+\sqrt{205})}{4}}\right).
\]
For this totally real tame quaternionic field of discriminant
$205^6$ from NN-A018,
\[
               \operatorname{ord}_{s=1/2}\zeta_N(s)=2.
\]
Its unique irreducible degree-two Artin factor has central order
exactly one. These statements require neither A nor M.

For the specific shifted ideal-norm series
$\zeta_{N,t}(s)=\sum_{n\ge1}a_N(n)(n+t)^{-s}$, continued as in
NN-A004, there are exactly two zero branches near $1/2$. Both
are strictly to the right for all sufficiently small real
$t>0$. More precisely, put
\[
 a_2=\frac{\zeta_N''(1/2)}2<0,\qquad
 A_N=\left(\frac{\zeta_N(3/2)}{2|a_2|}\right)^{1/2},\qquad
 d_N=-\frac{C_N A_N^2}{2}>0,
\]
where $C_N$ is the coefficient of NN-A018. Then
\begin{equation}\label{eq:analytic-martinet-two-right}
 s_\pm(t)=\frac12\pm i A_N\sqrt t+d_Nt+O(t^{3/2}),
 \qquad
 \Re s_\pm(t)-\frac12=d_Nt+O(t^2).
\end{equation}
In particular, their real velocities are
$\frac{d}{dt}\Re s_\pm(t)=d_N+O(t)>0$ on a sufficiently small
positive interval.
\end{corollary}

\begin{proof}
The field, discriminant and root-number input are those already
verified in NN-A018: Martinet, 1971, Propositions~III.1--III.2,
Theorem~IV.1 and Example~3, pp.~404--407, together with
Fr\"ohlich, 1974, Theorem~1, p.~82. The three quadratic central
values for discriminants $5,41,205$ are strictly positive by
the saved interval certificate of NN-A018. Hence the
factorization
\[
 \zeta_N(s)=\zeta(s)L(s,\chi_5)L(s,\chi_{41})
                         L(s,\chi_{205})L(s,\psi)^2
\]
and the negative symplectic root number imply
$m_N=2\operatorname{ord}_{1/2}L(s,\psi)\equiv2\pmod4$.
In particular $m_N\ge2$ and $a_{m_N}<0$; none of these steps
uses a Riemann hypothesis.

Apply the theorem at $\delta=1$. Since $r_1=8$, $r_2=0$ and
$\log D_N=6\log205$,
\[
 B_N(3/2)=\frac83+3\log205-4\log\pi+4\psi(3/4),
 \qquad
 B_N'(3/2)=-\frac{40}9+2\psi'(3/4).
\]
The reflection and duplication formulas for $\psi$, with
$\psi(1)=-\gamma_{\!E}$, give
\[
 \psi(3/4)=-\gamma_{\!E}+\frac\pi2-3\log2.
\]
For clarity, the other special value follows directly from
DLMF~5.15.1:
\[
 \psi'(3/4)=16\sum_{n\ge0}\frac1{(4n+3)^2}
                  =\pi^2-8\mathcal G,
 \qquad
 \mathcal G=\sum_{n\ge0}\frac{(-1)^n}{(2n+1)^2}.
\]
Indeed the sum of the two odd residue-class sums is
$\pi^2/8$ and their difference is $\mathcal G$.
Consequently
\begin{equation}\label{eq:analytic-martinet-U-exact}
 U_N(1)=\frac{32}9+\frac32\log205-2\log\pi
       -2\gamma_{\!E}+\pi-6\log2-\pi^2+8\mathcal G.
\end{equation}

Only coarse rational estimates are needed:
\begin{equation}\label{eq:analytic-martinet-coarse}
 \frac{157}{50}<\pi<\frac{22}7,\quad
 \log205<\frac{16}3,\quad \log\pi>\frac{57}{50},\quad
 \gamma_{\!E}>\frac12,\quad \log2>\frac{69}{100},\quad
 \mathcal G<\frac{23}{25}.
\end{equation}
Here are elementary finite checks, included so that no
floating-point constant is required. The Machin identity
$\pi=16\arctan(1/5)-4\arctan(1/239)$ follows by the tangent
addition law and the angle ranges. Alternating-series bounds
give
\[
 \begin{split}
 16\left(\frac15-\frac1{3\cdot5^3}\right)-\frac4{239}
                         &>\frac{157}{50},\\
 16\left(\frac15-\frac1{3\cdot5^3}
                         +\frac1{5\cdot5^5}\right)
 -4\left(\frac1{239}-\frac1{3\cdot239^3}\right)
                         &<\frac{22}7.
 \end{split}
\]
The positive series $\log x=2\sum_{j\ge0}y^{2j+1}/(2j+1)$,
$y=(x-1)/(x+1)$ for $x>1$, gives
\[
 \log\pi>\log(157/50)
       >2\sum_{j=0}^3\frac{(107/207)^{2j+1}}{2j+1}
       >\frac{57}{50},\qquad
 \log2>2\left(\frac13+\frac1{81}\right)
       =\frac{56}{81}>\frac{69}{100}.
\]
For the two upper logarithm bounds used here, the positive
exponential series gives the exact rational comparisons
\[
 \begin{split}
 \sum_{j=0}^{16}\frac{(16/3)^j}{j!}-205
       &=\frac{58220464777647308}{27485885585032875}>0,\\
 \sum_{j=0}^{8}\frac{(39/20)^j}{j!}-7
       &=\frac{3130848254169}{114688000000000}>0.
 \end{split}
\]
The first proves $\log205<16/3$. The elementary integral
comparison $\gamma_{\!E}>H_6-\log7$, the second comparison and
$H_6=49/20$ give $\gamma_{\!E}>1/2$. Finally, the alternating
Catalan series gives
\[
 \mathcal G<\sum_{j=0}^6\frac{(-1)^j}{(2j+1)^2}
       =\frac{23}{25}-\frac{3085982}{2029052025}
       <\frac{23}{25}.
\]
This proves every inequality in
\eqref{eq:analytic-martinet-coarse}. Substitution in
\eqref{eq:analytic-martinet-U-exact} now gives
\[
 \begin{split}
 U_N(1)&<\frac{32}9+8-\frac{57}{25}-1+\frac{22}7
            -\frac{207}{50}-\frac{24649}{2500}
            +\frac{184}{25}\\
       &=\frac{752663}{157500}
         =5-\frac{34837}{157500}<5.
 \end{split}
\]
The theorem yields $m_N<5$. Together with
$m_N\equiv2\pmod4$, this forces $m_N=2$ and the Artin order
one.

NN-A018 already proves $C_N<-3.5197<0$ using its saved
certificate. Its full local formula, now with $m=2$, gives
the first expansion in \eqref{eq:analytic-martinet-two-right}.
There are exactly two holomorphic sections after $t=u^2$.
Real coefficients and uniqueness identify these sections both
by conjugation on real $u$ and by $u\mapsto-u$; thus the odd
coefficients of either section are imaginary and its even
coefficients are real. The real part consequently has a
convergent power series in $t$, with first coefficient
$d_N>0$. This proves the improved $O(t^2)$ remainder, permits
termwise differentiation, and establishes the strict side and
velocity assertions.
\end{proof}

\paragraph{Classification and source boundary.}
The quaternionic root-number and central-existence results are
classical. Sami Omar, \emph{On Artin $L$-Functions for Octic
Quaternion Fields}, \emph{Experimental Mathematics}~10 (2001),
237--245, DOI~10.1080/10586458.2001.10504445, studies exact
central orders in these fields. The present proof does not
claim to be the first determination of the order of this field,
nor does it import a table entry that has not been inspected.
Its order bound is fully proved above and uses the existing
NN-A018 certificate only for the quadratic nonvanishing and
the second-jet sign. The comparison with an exact Omar table
entry remains a source-location audit, not a missing hypothesis
of the deduction.

\paragraph{AG and representation translation.}
The symplectic irreducible representation of the actual
$Q_8$-extension has negative functional-equation sign and,
by the proved upper bound, a length-one central analytic zero.
Its multiplicity two in the regular representation produces
the local Dedekind Artin algebra $\mathbb C[z]/(z^2)$. For the
specified shifted deformation, the base change $t=u^2$ gives
two sections with a common positive second real normal
coefficient $d_N$. Their discriminant is
$-4A_N^2t+O(t^2)$, so the cover has simple ramification at
$t=0$, is etale off that point on a sufficiently small disc,
and has the transposition monodromy. This is the concrete
finite-free discriminant criterion of SGA~1, Expos\'e~I,
section~4, Proposition~10. Positivity of the rational test is
a statement about conjugate analytic zero orbits, not
Frobenius positivity or a claimed motive for the shifted
series. The sign of the second normal jet retains the analytic
unit data that the length and the reduced support alone do not
record.
